\begin{abstractpt}
As ferramentas assistidas por IA de uso geral, como o ChatGPT, têm recentemente ganhado muita atenção da mídia e do público em geral, o que levantou questões sobre em quais tarefas podemos aplicar tal ferramenta.Um bom design de código é essencial para o desenvolvimento ágil de software para mantê-lo pronto para mudanças. Nesse contexto, identificar qual padrão de projeto pode ser apropriado para um determinado cenário pode ser considerada uma habilidade avançada que requer um alto grau de abstração e um bom conhecimento de orientação a objetos. Este artigo tem como objetivo realizar um estudo exploratório investigando a eficácia de uma ferramenta assistida por IA para auxiliar desenvolvedores na escolha de um padrão de projeto para resolver cenários de projeto. Para atingir esse objetivo, reunimos 56 perguntas existentes usadas por professores e concursos públicos que fornecem um contexto concreto e perguntam qual padrão de design seria adequado. Enviámos estas questões ao ChatGPT e analisámos as respostas. Constatamos que 93\% das perguntas foram respondidas corretamente com um bom nível de detalhamento, demonstrando o potencial dessa ferramenta como um recurso valioso para ajudar desenvolvedores a aplicar padrões de projeto e tomar decisões de projeto
\end{abstractpt}
